\section{Conclusion}\label{sec:conc}
This study set out to build a lesion-aware system for the early detection of diabetic retinopathy and to test it honestly on external data. The lesion-aware system detected more mild cases than the global-pooling baseline, significantly so on the internal and APTOS sets, but it did so by operating at a lower specificity, and its discrimination, measured by AUC, was slightly worse. On the external sets from India and China its sensitivity stayed above 0.90, yet specificity and the positive predictive value at a 10\% prevalence remained far below clinical targets. Neither the built-in attention nor Grad-CAM pointed to annotated lesions more often than a random point, so neither can be offered to graders as an explanation. The quality gate itself was accurate, but it did not change accuracy on the images it accepted, and it rejected severe and proliferative images much more often than healthy ones.

The practical significance lies mainly in the audit. Public test sets contain duplicated images with conflicting grades, including a leak across the official IDRiD split, so results on them must be preceded by de-duplication. At the 224-pixel resolution used by most models, half of all microaneurysms are smaller than one pixel; a model at this resolution cannot see the lesion that defines the earliest stage, and improvements of the classifier head cannot compensate. Because poor image quality is concentrated in proliferative disease, a quality gate in a screening pathway must route rejected images to referral and never report them as negative.

The study has clear limitations. It was run on two CPU cores: the planned retinal foundation model was replaced by a partially fine-tuned EfficientNet-B0, only one training seed and five epochs were possible, and the ablations of the loss and of a higher input resolution were not run. Messidor-2 with its adjudicated grades and the graded part of DDR were not available, the EyePACS reference standard is noisy at the grade~0/1 boundary, and no clinical data or local images from Kazakhstan were included. The results are therefore a lower bound for the approach rather than a test of its full form.

Future work follows from these findings: training at 448 pixels or higher, or with lesion-centred crops, to keep microaneurysms visible; replacing the backbone with a retinal foundation model while excluding its EyePACS pretraining images from internal testing; running five or more seeds for every ablation; evaluating on Messidor-2 and on a local cohort under ethics approval; and supervising or calibrating the attention maps with the available lesion masks so that their localisation can exceed chance.
